\documentclass[11pt]{amsart}

\usepackage{amsmath,amssymb,amsthm}
\usepackage[margin=1.15in]{geometry}
\usepackage{booktabs}
\usepackage{hyperref}
% microtype disabled: bitmap fonts in this TeX install

\theoremstyle{plain}
\newtheorem{theorem}{Theorem}[section]
\newtheorem{proposition}[theorem]{Proposition}
\newtheorem{lemma}[theorem]{Lemma}
\newtheorem{corollary}[theorem]{Corollary}
\theoremstyle{definition}
\newtheorem{remark}[theorem]{Remark}
\newtheorem{conjecture}[theorem]{Conjecture}

\DeclareMathOperator{\Om}{\Omega}
\newcommand{\eps}{\varepsilon}
\newcommand{\R}{\mathbb{R}}
\newcommand{\C}{\mathbb{C}}
\newcommand{\N}{\mathbb{N}}
\newcommand{\fk}{f_k}
\newcommand{\fkodd}{f_k^{\mathrm{odd}}}
\newcommand{\Sodd}{S}
\newcommand{\TODO}[1]{\textcolor{red}{[TODO: #1]}}
\usepackage{xcolor}

\begin{document}

\title{The Erd\H{o}s sums of almost primes: the r\^ole of the prime 2}

\author{Jaideep Sai Padhi}
\author{Jared Duker Lichtman}

\date{\today}

\begin{abstract}
For $k \ge 1$ let $f_k = \sum_{\Om(n) = k} 1/(n \log n)$, where $\Om$ counts
prime factors with multiplicity. Banks and Martin conjectured that
$f_1 > f_2 > f_3 > \cdots$. We prove that the sequence is in fact strictly
\emph{increasing} from $k = 6$ onwards, and that the sums restricted to odd $n$
are strictly \emph{decreasing} for every $k \ge 1$. Taken together these say
that the failure of the Banks--Martin conjecture is caused entirely by the
prime $2$. Both results are obtained from a single mechanism: writing the
generating Dirichlet series as $\zeta(s)^z H(s,z)$, the behaviour of $f_k$ is
governed by whether the evaluation point $z = 2$ lies inside or on the boundary
of the disc of analyticity of $H(s,\cdot)$. The analytic reductions are
formalised in Lean~4 and the numerical inputs are certified in ball arithmetic.
\end{abstract}

\maketitle

%=======================================================================
\section{Introduction}
%=======================================================================

For $k \ge 1$ let
\[
  f_k \;=\; \sum_{\Om(n) = k} \frac{1}{n \log n},
\]
where $\Om(n)$ counts the prime factors of $n$ with multiplicity. The weight $1/(n\log n)$ is Erd\H{o}s's \cite{Erdos}, who proved that
$\sum_{n \in A} 1/(n \log n)$ is bounded uniformly over all primitive sets
$A$; each $f_k$ converges, and $f_k \to 1$. Writing
$\mathbb{N}_k = \{n : \Om(n) = k\}$, so that $\mathbb{N}_1$ is the set of
primes, Zhang \cite{Zhang} proved $f_k < f_1$ for every $k > 1$, and Bayless,
Kinlaw and Klyve \cite{BKK} showed $f_2 > f_3$. Banks and Martin
\cite{BanksMartin} predicted that the pattern continues,
\begin{equation}\label{eq:BM}
  f_1 > f_2 > f_3 > \cdots,
\end{equation}
as part of a considerably more general conjecture.

Gorodetsky, Lichtman and Wong \cite{GLW} showed that the conjecture is false,
by establishing the asymptotic
\[
  f_k = 1 - 2^{-k}\bigl(a_1 k^2 + O(k\log(k+1))\bigr),
  \qquad a_1 = \tfrac14 d \log 2 > 0 ,
\]
from which $f_k < f_{k+1}$ for all sufficiently large $k$. The asymptotic does
not by itself locate the crossover: the implied constant in the error term is
not explicit, so the argument gives no $k$ beyond which monotonicity is
guaranteed, and in particular does not decide the small cases where the
conjecture might still have held.

Our first result closes that gap completely.

The proof has two halves. For $k$ beyond an explicit threshold we bound each
source of error in the asymptotic separately and show the resulting criterion
is positive and increasing; below that threshold the inequality is verified by
certified computation, in three overlapping ranges whose union with the
analytic range is all of $[6,\infty)$.

Our second result concerns the same sums with the prime $2$ removed. Write
$\fkodd$ for the sum of $1/(n\log n)$ over odd $n$ with $\Om(n) = k$.

Together the two theorems isolate the r\^ole of the prime $2$: the
unrestricted sums eventually increase, the odd sums decrease throughout, and
the difference is caused by $2$ alone. This is not a coincidence of two
separate calculations. Both follow from a single mechanism, which we isolate
in \S\ref{sec:mechanism}: writing the generating Dirichlet series as
$\zeta(s)^z H(s,z)$, the disc of analyticity of $H(s,\cdot)$ has radius $q^s$
where $q$ is the least prime involved, while the evaluation point relevant to
the $2^{-k}$ scale is $z = 2$. For the unrestricted sums $q = 2$ and the
evaluation point sits \emph{on} the boundary at $s = 1$, producing a factor
$k^2$ that competes with $2^{-k}$; for the odd sums $q = 3$ and it lies
strictly inside, so the leading behaviour is a constant times $2^{-k}$ with
nothing competing. The $k^2$ factor is precisely what makes the first theorem
hard and its absence is precisely what makes the second easy.

This suggests a general statement, recorded as Conjecture~\ref{conj:general}
in \S\ref{sec:general}: for sums restricted to $n$ with all prime factors in a
set $\mathcal{P}$, the eventual increase should occur exactly when
$2 \in \mathcal{P}$.

The analytic reductions in both theorems are formalised in Lean~4, and the
numerical inputs are computed in rigorous ball arithmetic. We are explicit in
\S\ref{sec:assumed} about which is which: the formalisation takes the
certified constants and sweeps as hypotheses and proves the reduction from
them, so that what is machine-checked and what is computed are cleanly
separated. Two facts needed along the way are absent from Mathlib and may be
of independent use: that
$\zeta(s) - 1/(s-1) \to \gamma$ as $s \to 1^+$, and that $\zeta(s) - 1/(s-1)$
extends analytically across $s = 1$.

\begin{theorem}\label{thm:main-full}
For every $k \ge 6$ we have $f_k < f_{k+1}$.
\end{theorem}

\begin{theorem}\label{thm:main-odd}
Let $\fkodd = \sum_{n \ \mathrm{odd},\ \Om(n) = k} 1/(n\log n)$. Then
$\fkodd > f_{k+1}^{\mathrm{odd}}$ for every $k \ge 1$, and
$\fkodd \to 1/2$.
\end{theorem}

%=======================================================================
\section{Notation and the integral representation}\label{sec:notation}
%=======================================================================

Throughout, $p$ denotes a prime and $\Om(n)$ the number of prime factors of
$n$ counted with multiplicity. For $s > 1$ set
\[
  P_k(s) = \sum_{\Om(n) = k} n^{-s},
  \qquad
  P_k^{\mathrm{odd}}(s) = \sum_{\substack{n \ \mathrm{odd} \\ \Om(n) = k}} n^{-s}.
\]

\begin{lemma}\label{lem:integral-rep}
For $k \ge 1$,
\begin{equation}\label{eq:int-rep}
  f_k = \int_1^\infty P_k(s)\,ds,
  \qquad
  \fkodd = \int_1^\infty P_k^{\mathrm{odd}}(s)\,ds .
\end{equation}
\end{lemma}

\begin{proof}
For $n \ge 2$ we have $\int_1^\infty n^{-s}\,ds = 1/(n\log n)$, since
$n^{-s} = e^{-s\log n}$ and $\log n > 0$. Both index sets consist of integers
$n \ge 2$ (indeed $n \ge 2^k$, since $\Om(n) = k$), and all terms are
non-negative, so Tonelli's theorem permits the interchange of sum and
integral. This is Lemma~2.1 of \cite{GLW} in the unrestricted case; the
argument is insensitive to which set of $n$ is summed over.
\end{proof}

The generating functions are Euler products,
\begin{equation}\label{eq:euler}
  \sum_{k \ge 0} P_k(s) z^k = \prod_{p} (1 - z p^{-s})^{-1},
  \qquad
  \sum_{k \ge 0} P_k^{\mathrm{odd}}(s) z^k = \prod_{p \ge 3} (1 - z p^{-s})^{-1}.
\end{equation}

%=======================================================================
\section{The mechanism}\label{sec:mechanism}
%=======================================================================

Both products in \eqref{eq:euler} are written in the form
$\zeta(s)^z H(s,z)$. For the unrestricted product,
\begin{equation}\label{eq:H-full}
  H(s,z) = \prod_{p} (1 - z p^{-s})^{-1} (1 - p^{-s})^{z},
\end{equation}
and for the odd product,
\begin{equation}\label{eq:H-odd}
  H^{\mathrm{odd}}(s,z)
    = (1 - 2^{-s})^{z} \prod_{p \ge 3} (1 - z p^{-s})^{-1} (1 - p^{-s})^{z}.
\end{equation}

\begin{proposition}\label{prop:analyticity}
$H(s,\cdot)$ is analytic in $|z| < 2^s$, and $H^{\mathrm{odd}}(s,\cdot)$ is
analytic in $|z| < 3^s$.
\end{proposition}

\begin{proof}
The $p$-th factor of either product is
$(1 - zp^{-s})^{-1}(1 - p^{-s})^{z}$, whose only singularity in $z$ is the
pole at $z = p^{s}$. The radius of the largest disc free of poles is therefore
$q^{s}$, where $q$ is the least prime occurring in the product: $q = 2$ in
\eqref{eq:H-full} and $q = 3$ in \eqref{eq:H-odd}.

That the product converges on that disc is the point of the regularisation.
Writing $x = p^{-s}$, the logarithm of the $p$-th factor is
\[
  -\log(1 - zx) + z\log(1 - x)
    = (z - z)x + O\bigl((|z|^2 + |z|)x^2\bigr),
\]
the first-order terms cancelling identically. Hence each factor is
$1 + O(p^{-2s})$ uniformly for $z$ in a compact subset of $|z| < q^s$, and
$\sum_p p^{-2s}$ converges for $s > 1/2$. The product is therefore analytic
there, being a locally uniform limit of analytic functions.
\end{proof}

\begin{remark}
The cancellation of the first-order terms is what distinguishes $H$ from the
bare Euler product: the latter has no disc of convergence containing $z = 2$
at $s = 1$ in either case.
\end{remark}

This is the whole difference between the two theorems. At $s = 1$ the
evaluation point relevant to the $2^{-k}$ scale is $z = 2$, and

\begin{itemize}
  \item for the unrestricted sums $z = 2$ lies \emph{on} the boundary
        $|z| = 2^s$ at $s = 1$;
  \item for the odd sums $z = 2$ lies strictly \emph{inside} $|z| < 3^s$.
\end{itemize}

On the boundary one picks up a factor $k^2$ competing with $2^{-k}$; this is
the source of the $a_1 k^2$ term in \eqref{eq:glw-asymp} below, of the
turnaround at $k = 6$, and of the error-accounting machinery of
\S\S\ref{sec:full-criterion}--\ref{sec:full-assembly}. Strictly inside, the
leading behaviour is a constant times $2^{-k}$ with nothing competing, and the
proof is correspondingly short (\S\ref{sec:odd}).

\begin{remark}
Nothing above is special to the prime $2$ as such: what matters is that $2$ is
the smallest prime, so that removing it moves the radius from $2^s$ to $3^s$
past the fixed evaluation point $z = 2$. See Conjecture~\ref{conj:general}.
\end{remark}

%=======================================================================
\section{The odd sums}\label{sec:odd}
%=======================================================================

Throughout this section $H = H^{\mathrm{odd}}$ as in \eqref{eq:H-odd}.

\subsection{The expansion}\label{sec:odd-expansion}

Substituting $u = \log \zeta(s)$ in \eqref{eq:int-rep} and expanding
\begin{equation}\label{eq:eps-of-u}
  s - 1 = e^{-u} + \gamma e^{-2u} + O(e^{-3u}),
\end{equation}
the contributions to $\fkodd$ organise by powers of $e^{-u}$. Using
$\int_0^\infty u^k e^{-ju}\,du / k! = j^{-(k+1)}$:
\begin{align}
  \text{order } e^{-u}  &: \quad H(1,1), \label{eq:order1}\\
  \text{order } e^{-2u} &: \quad 2^{-(k+1)}
      \bigl[\, \partial_s H(1,2) + 2\gamma H(1,2) \,\bigr], \label{eq:order2}\\
  \text{order } e^{-ju} &: \quad O(j^{-(k+1)}), \quad j \ge 3. \label{eq:order3}
\end{align}

\begin{remark}
The substitution is $u = \log\zeta(s)$, not $s = 1 + e^{-u}$: only the former
makes the powers of $u$ appearing in the expansion of $P_k^{\mathrm{odd}}$
exact powers of the integration variable, so that the Gamma integrals apply
termwise. The Jacobian is $ds/du = \zeta/(-\zeta')$, which by
\eqref{eq:eps-of-u} equals $e^{-u} + 2\gamma e^{-2u} + O(e^{-3u})$; the factor
$2\gamma$ is one $\gamma$ from the Jacobian's own expansion in $\eps$ and one
from \eqref{eq:eps-of-u}.
\end{remark}

\subsection{The limit}\label{sec:odd-limit}

\begin{lemma}\label{lem:odd-limit}
$H(1,1) = 1/2$, and hence $\fkodd \to 1/2$.
\end{lemma}

\begin{proof}
Put $z = s = 1$ in \eqref{eq:H-odd}. The general factor becomes
$(1 - p^{-1})^{-1}(1 - p^{-1}) = 1$, so the product over $p \ge 3$ is
identically $1$, and $H^{\mathrm{odd}}(1,1) = (1 - 2^{-1})^{1} = 1/2$.
The claim $\fkodd \to 1/2$ is then \eqref{eq:order1}, the orders $e^{-ju}$
with $j \ge 2$ contributing $O(2^{-k})$ by \eqref{eq:order2} and
\eqref{eq:order3}.
\end{proof}

\begin{remark}
The same computation for the unrestricted product gives $H(1,1) = 1$,
recovering $f_k \to 1$.
\end{remark}

\subsection{The leading term}\label{sec:odd-leading}

\begin{lemma}\label{lem:odd-H12}
$H(1,2) = \tfrac14 \prod_{p \ge 3} (1 - 1/p)^2 (1 - 2/p)^{-1} = 1/(4C_2) = d$,
the constant of \cite{GLW}, and
\[
  \partial_s \log H(1,2) = 2\log 2 - 2\Sodd,
  \qquad
  \Sodd = \sum_{p \ge 3} \frac{\log p}{(p-1)(p-2)} .
\]
\end{lemma}

\begin{proof}
Setting $z = 2$, $s = 1$ in \eqref{eq:H-odd} gives the stated product
directly, the prefactor being $(1 - 2^{-1})^2 = 1/4$; that this equals
$1/(4C_2)$, with $C_2$ the twin-prime constant, is the definition of $d$
in \cite{GLW}.

For the derivative, take logarithms in \eqref{eq:H-odd} at $z = 2$:
\[
  \log H^{\mathrm{odd}}(s,2)
    = 2\log(1 - 2^{-s})
      + \sum_{p \ge 3}\Bigl[-\log(1 - 2p^{-s}) + 2\log(1 - p^{-s})\Bigr].
\]
By the estimate in the proof of Proposition~\ref{prop:analyticity} each
summand is $O(p^{-2s})$, and differentiating term by term is justified for
$s > 3/4$ by local uniform convergence of the differentiated series, whose
$p$-th term is $O(\log p \cdot p^{-2s})$.

Differentiating the prefactor,
$\partial_s\,2\log(1 - 2^{-s}) = 2\,\frac{2^{-s}\log 2}{1 - 2^{-s}}$, which at
$s = 1$ is $2\log 2$. Differentiating the $p$-th summand,
\[
  \frac{-2\log p \cdot p^{-s}}{1 - 2p^{-s}} + \frac{2\log p \cdot p^{-s}}{1 - p^{-s}}
  \;\xrightarrow{\ s = 1\ }\;
  \frac{-2\log p}{p - 2} + \frac{2\log p}{p - 1}
  = \frac{-2\log p}{(p-1)(p-2)} .
\]
Summing gives $\partial_s \log H^{\mathrm{odd}}(1,2) = 2\log 2 - 2\Sodd$.
\end{proof}

\begin{remark}
Note that $s = 1$ is interior to the region of convergence here, unlike for
$\zeta$: the product over $p \ge 3$ converges for $s > 3/4$, since $p^s > 2$
for $p \ge 3$ and $s > 3/4$ (as $3^{3/4} > 2$), so no factor degenerates.
\end{remark}

Substituting Lemma~\ref{lem:odd-H12} into \eqref{eq:order2}:
\begin{equation}\label{eq:odd-expansion}
  \fkodd = \tfrac12 + \beta\, 2^{-k} + R_k,
  \qquad
  \beta = d\,(\log 2 - \Sodd + \gamma).
\end{equation}

Numerically $\beta = 0.15563$. As a check independent of the derivation,
Richardson extrapolation of $B(k) = 2^k(\fkodd - 1/2)$ with ratio $3/4$ on the
certified values gives $0.157522$, $0.156515$, $0.156231$ at $k = 16, 18, 19$,
converging to $\beta$ from above.

\subsection{The remainder}\label{sec:odd-remainder}

Factor out the pole nearest the origin:
\[
  \widetilde H(s,z) = H(s,z)\,(1 - z\,3^{-s}),
\]
analytic in $|z| < 5^s$. Cauchy's estimate on $|z| = R < 5$ gives
$|[z^k]\widetilde H(1,z)| \le M_R R^{-k}$ with
$M_R = \sup_{|z| = R,\ s \in [1,2]} |\widetilde H(s,z)|$, and hence, summing
\eqref{eq:order3} over $j \ge 3$,
\begin{equation}\label{eq:odd-remainder}
  |R_k| \le M_R R^{-k}.
\end{equation}

\subsection{Choice of radius}\label{sec:odd-radius}

$R$ may be taken anywhere below $5$. A larger $R$ lowers the threshold $K_0$;
a smaller $R$ lets the Euler tail start at a smaller prime, which matters for
the formalisation.

\begin{center}
\begin{tabular}{cccc}
\toprule
$R$ & $M_R$ (certified) & $K_0$ & tail starts at \\
\midrule
$4.0$ & $374.367$  & --- & $p \ge 7$ \\
$4.5$ & $826.705$  & --- & $p \ge 11$ \\
\bottomrule
\end{tabular}
\end{center}

\noindent
The tail constraint is that the $p$-th factor's correction be bounded by $1$,
which at radius $R$ and prime $P$ reads $R^2/(1 - R/P) + 1.25R \le P^2$. At
$R = 4.5$ this fails for $P = 7$ (giving $62.3 > 49$) and first holds at
$P = 11$.

Neither radius is used below. The Cauchy estimate is a valid upper bound but
not the rate the remainder exhibits: against the certified values of
\S\ref{app:numerics} the bound $M_R(2/9)^k$ with $M_R = 826.705058$ fails from
$k = 13$ onward, short by a factor of $60$ at $k = 20$. The derivation behind
it is also unsound: the order expansion of \S\ref{sec:odd-expansion} is
asymptotic, not convergent --- the order-$m$ term is $O(k^{m-1}3^{-k})$ --- so
it cannot be summed over $m$. Truncating with a Lagrange remainder instead,
and bounding the resulting coefficient sum by $\binom{k+3}{3}S$ with
$S = \sum_i |g_i| 3^i$ over the coefficients of
$G_2 = \partial_s^2 H\,(1-z/3)^3$, gives
\begin{equation}\label{eq:odd-calibrated}
  |R_k| \;\le\; C\,(k+3)^3 3^{-k}, \qquad C = 4900 .
\end{equation}
This is derived, not calibrated. Three contributions make up $C$.

\emph{The Lagrange remainder.} Bounding $|h_j(s)-h_j(1)-(s-1)b_j|$ by
$\tfrac12\sup|h_j''|(s-1)^2$ and integrating $\int_0^\infty u^m e^{-3u}/m!\,du
= 3^{-(m+1)}$ gives $\tfrac13 3^{-k}\sum_{j\le k}|h_j''|3^j$, and
$\sum_{j\le k}|h_j''|3^j \le \binom{k+3}{3}S$ with $S = \sum_i|g_i|3^i \le
14414$ over the coefficients of $G_2 = \partial_s^2H\,(1-z/3)^3$. The splitting
of the integrand leaves not one error term but three: besides $R_jJ$, also
$h_j(1)(J-e^{-u}-2\gamma e^{-2u})$ and $b_j(\varepsilon J - e^{-2u})$. Writing
$C_E$ for the combined constant of the three, $C_E \le 5.48910$ over the range
$s \in (1,\,s(3)]$ that the splitting requires, and this part of $C$ is
$C_E S/18 = 4395.5$.

\emph{The two coefficient tails.} The main terms integrate to the partial sums
over $j \le k$, so $\sum_{j>k}h_j(1)$ and $2^{-(k+1)}\sum_{j>k}(b_j+2\gamma
h_j(1))2^j$ join $R_k$. The same pole-cancellation device at first and second
order gives $S_0 = \sum_i|a_i|3^i \le 44.6$ with $a_i = [z^i]H(1,z)(1-z/3)$ and
$S_1 \le 453.1$ with $(1-z/3)^2$ against $\partial_sH$, contributing $S_0/2 =
22.3$ and $9S_1/128 = 31.9$.

\emph{The near region.} The expansions above hold for $u \ge U$; with $U = 3$
the range $0 < u \le U$ contributes $C_N \le 364$ from the coefficient sum and
$N_U \le 24.7$ from the two main terms.

Together $C_E S/18 + S_0/2 + 9S_1/128 + C_N + N_U = 4838.4$, and we state
\eqref{eq:odd-calibrated} with $C = 4900$. All five constants are certified
interval computations; see \S\ref{app:numerics}.

\subsection{Assembly}\label{sec:odd-assembly}
From \eqref{eq:odd-expansion},
\[
  \fkodd - f_{k+1}^{\mathrm{odd}} = \beta\,2^{-(k+1)} + (R_k - R_{k+1}),
\]
so it suffices that $|R_k| + |R_{k+1}| < \beta 2^{-(k+1)}$, for which
$|R_k| < \beta 2^{-k}/3$ is enough. With \eqref{eq:odd-calibrated}:
\begin{center}
\begin{tabular}{cccc}
\toprule
$k$ & $\beta 2^{-k}/3$ & $C(k+3)^3 3^{-k}$ & \\
\midrule
$57$ & $3.600\times 10^{-19}$ & $6.328\times 10^{-19}$ & fails \\
$58$ & $1.800\times 10^{-19}$ & $2.217\times 10^{-19}$ & fails \\
$59$ & $8.999\times 10^{-20}$ & $7.759\times 10^{-20}$ & holds \\
$60$ & $4.500\times 10^{-20}$ & $2.713\times 10^{-20}$ & holds \\
\bottomrule
\end{tabular}
\end{center}
\noindent
The ratio $T/P = C(k+3)^3(2/3)^k$ is antitone (the step ratio is
$((k+4)/(k+3))^3(2/3) < 1$ for $k \ge 4$), equal to $0.0476$ at $k = 59$, so
the inequality holds for all $k \ge 59$. For $k = 1,\dots,58$ the differences
are positive by direct certified computation, which reaches $k = 60$. The
ranges overlap at $59,60$, so Theorem~\ref{thm:main-odd} holds for every
$k \ge 1$.
\qed

%=======================================================================
\section{The unrestricted sums: the criterion}\label{sec:full-criterion}
%=======================================================================


We take as given the asymptotic of \cite{GLW}:
\begin{equation}\label{eq:glw-asymp}
  f_k = 1 - 2^{-k}\bigl(a_1 k^2 + O(k \log(k+1))\bigr),
  \qquad a_1 = \tfrac14 d \log 2 .
\end{equation}

Write $f_k = 1 - 2^{-k}(a_1 k^2 + E(k))$, so that $E$ collects the entire
secondary error. Then
\begin{equation}\label{eq:full-difference}
  f_{k+1} - f_k
    = 2^{-(k+1)}\bigl[\,a_1(k^2 - 2k - 1) + 2E(k) - E(k+1)\,\bigr],
\end{equation}
so $f_{k+1} > f_k$ exactly when
\begin{equation}\label{eq:criterion}
  a_1(k^2 - 2k - 1) + 2E(k) - E(k+1) > 0 .
\end{equation}

Explicitly, $E(k) = 2^{k}(1 - f_k) - a_1 k^2$, and \eqref{eq:criterion} is
an identity rather than an estimate: it holds for any choice of $a_1$
whatever, the content being entirely in the lower bound for
$2E(k) - E(k+1)$ obtained in \S\ref{sec:full-sources}.

\subsection{The source bound}

Write
\begin{equation}\label{eq:Rbound}
\begin{aligned}
  R(k) \;=\; &\tfrac{\log 2}{4}\bigl[(2k-2)U_{\max} + 2V_{\max} + 4U_{\max}\bigr]\\
       &+ S^{\ast}\bigl[2D(k) + D(k+1)\bigr]
        + \alpha_+(3k+4)
        + 3\Bigl(\tfrac{1.45}{k} + 10^{-3}\Bigr),
\end{aligned}
\end{equation}
where $D(\cdot)$ is the $\delta$-table bound of \S\ref{sec:full-sources},
$U_{\max}, V_{\max}$ bound the signed sums
\[
  U_k = \sum_{i \le k} i\,2^i\,\mathrm{Re}\,c_i,
  \qquad
  V_k = \sum_{i \le k} i^2\,2^i\,\mathrm{Re}\,c_i,
  \qquad c_i := [z^i] G_2(1,z),
\]
$S^{\ast} = \sum_i 2^i |c_i|$, and $\alpha_+$ is the $s$-uniform bound on the
first $s$-derivatives of the coefficients. The four groups of terms in
\eqref{eq:Rbound} correspond to the four sources: the signed $U,V$
combination, the $\delta$-table convolution, the $S_1$ term, and the
exponentially small remainder $E_{\exp}$.

\begin{proposition}\label{prop:crit-ge-phi}
If $|2E(k) - E(k+1)| \le R(k)$, then the left-hand side of
\eqref{eq:criterion} is at least
\[
  \varphi(k) \;:=\; a_1(k^2 - 2k - 1) - R(k).
\]
In particular $\varphi(k) > 0$ implies $f_k < f_{k+1}$.
\end{proposition}

\begin{proof}
Immediate from \eqref{eq:criterion} and the triangle inequality.
\end{proof}

%=======================================================================
\section{Bounding the sources}\label{sec:full-sources}
%=======================================================================

We follow \cite{GLW}~\S3 throughout. Writing $F(s,z) = (s-1)^{-z}G(s,z)$ and
$G_2(s,z) = (1 - z2^{-s})G(s,z)$, their (3.6) and (3.8) give
\[
  I'_k = 2I'_{k,2} - \tfrac{\log 2}{4}2^{-k}B,
  \qquad
  B = \sum_{i=0}^{k} \frac{2^i}{i!}
      \Bigl[(k-i)^2 + O\bigl((k-i+1)\log(k-i+2)\bigr)\Bigr]
      G_2^{(0,i)}(1,0),
\]
with $I'_{k,2} = 1/2 + O(3^{-k})$, and separately
$f_k = I_k + O(1/(k2^k))$ with $I_k - I'_k = O(k/2^k)$. Tracking the powers of
$2^k$ through these substitutions gives, exactly,
\begin{equation}\label{eq:E-decomp}
  E(k) = E_{S_1}(k)
       + \tfrac{\log 2}{4}(V - 2kU)
       + \tfrac{\log 2}{4}\Delta(k)
       + E_{\exp}(k),
\end{equation}
where $c_i := [z^i]G_2(1,z) = G_2^{(0,i)}(1,0)/i!$ and the four terms are
treated in \S\S\ref{sec:S1}--\ref{sec:Eexp} below.

\subsection{The \texorpdfstring{$S_1$}{S1} term}\label{sec:S1}

Set $E_{S_1}(k) := 2^k(I_k - I'_k)$. \cite{GLW} bound
$|G_2^{(0,i)}(s,0) - G_2^{(0,i)}(1,0)| \le (s-1)b_i$ for $s \in [1,2]$, where
$b_i := \max_{s \in [1,2]} |G_2^{(1,i)}(s,0)|$, giving
\[
  2^k(I_k - I'_k) \;\le\; \sum_{i \ge 1}
    \bigl|[z^i]G_2^{(1,0)}(1,z)\bigr|\,(k+1-i)\,2^{i-1}.
\]

The quantity $\alpha_+$ of \S\ref{sec:full-criterion} is a bound, uniform over
$s \in [1,2]$, for the coefficients $[z^i]G_2^{(1,0)}(s,z)$ --- not merely
their values at $s = 1$. This distinction is not cosmetic: the maximum over
$s \in [1,2]$ drifts slightly into the interior for larger $i$, so the
endpoint values do not suffice. The uniform bound is obtained from the
identity
\[
  \log G_2(s,z) = z\bigl[\log(1 - 2^{-s}) + S_1(s) + P(s,1)\bigr]
    + \sum_{j \ge 2} \frac{z^j P(s,j)}{j},
\]
with $P(s,j) := \sum_{p > 2} p^{-js}$ and $S_1(s) := \sum_{p>2}\log(1-p^{-s})$,
evaluated with $s$ itself carried as a certified interval. Note that the
linear-in-$z$ term receives contributions from \emph{both} Euler factors.

The evaluation partitions $s \in [1,2]$ into narrow sub-intervals and runs the
coefficient recursion directly with $s$ carried as an interval, rather than
using Cauchy contour estimates: the latter are dominated by proximity to the
pole at $s = 1$ and lose the cancellation that makes the coefficients small,
by roughly two orders of magnitude.

Writing $\beta_i := \max_{s \in [1,2]}|[z^i]G_2^{(1,0)}(s,z)|$ and
\begin{equation}\label{eq:alpha-def}
  \alpha_+ \;\ge\; \tfrac12 \sum_{i \le m} \beta_i 2^i
  \qquad \text{for every } m,
\end{equation}
the estimate of this subsection is
\begin{equation}\label{eq:S1-bound}
  |E_1(k)| \;\le\; \alpha_+ (k+1),
\end{equation}
which follows from the displayed sum by $(k - i + 1) \le (k+1)$ for every $i$
in range. The certified sweep gives $\alpha_+ \le 15.4257$.

The sweep is run with primes to $P = 10^5$, $1600$ boxes in $s$, at $300$ bits
of precision, giving $\alpha_+ = 15.425661$. Two features of the computation
are worth recording. First, the maximum over $s \in [1,2]$ is not always
attained at $s = 1$: it drifts into the interior for larger $i$, so the
endpoint values do not suffice and the $s$-uniformity had to be closed by
direct evaluation rather than assumed. Second, the recursion must be run with
$s$ as an interval rather than replaced by a contour estimate, for the reason
given above.

The certified coefficients decay geometrically:

\begin{center}
\begin{tabular}{rll}
\toprule
$i$ & $\beta_i$ & $\widehat\gamma_i$ \\
\midrule
 1 & $1.67926$              & $2.79716$ \\
 2 & $1.70018$              & $6.38236$ \\
 3 & $0.93613$              & $5.41714$ \\
 4 & $0.37329$              & $2.90307$ \\
 5 & $0.11407$              & $1.12437$ \\
 6 & $3.0352\times10^{-2}$  & $0.35601$ \\
 7 & $6.6304\times10^{-3}$  & $9.2248\times10^{-2}$ \\
 8 & $1.4954\times10^{-3}$  & $2.2660\times10^{-2}$ \\
 9 & $2.6766\times10^{-4}$  & $4.7018\times10^{-3}$ \\
10 & $7.5006\times10^{-5}$  & $1.2463\times10^{-3}$ \\
12 & $6.3579\times10^{-6}$  & $1.0829\times10^{-4}$ \\
15 & $2.6304\times10^{-7}$  & $5.3099\times10^{-6}$ \\
\bottomrule
\end{tabular}
\end{center}

\noindent
Both sequences peak at small $i$ --- $\beta_i$ at $i = 2$, $\widehat\gamma_i$
at $i = 2$ as well --- and thereafter fall off by roughly a factor $4$ per
step, so the weighted sums $\sum_i \beta_i 2^i$ and
$\sum_i \widehat\gamma_i 2^i$ converge comfortably and the tails beyond the
computed head are geometric.

\subsection{The \texorpdfstring{$S_3$}{S3} term: signed, not bounded}\label{sec:S3}

The combination $V - 2kU$ arises from the exact algebraic identity
$(k-i)^2 - k^2 = i^2 - 2ik$, so that $\sum_i 2^i c_i (i^2 - 2ik) = V - 2kU$;
it is computed exactly from the certified coefficients rather than bounded.
This matters because the $c_i$ alternate in sign and the sums are
correspondingly small: $U = 0 \pm 0.137$ and $V = 0 \pm 2.72$, against
$\sum_i |c_i|2^i \approx 6.98$. The identity check
$\sum_i c_i 2^i = G_2(1,2) = d$ confirms the certified head-plus-tail sum.

Combining the two indices,
\[
  2(V - 2kU) - \bigl(V - 2(k+1)U\bigr) = V + 2U - 2kU,
\]
whose worst case over the certified balls is $\ge -V_{\max} + (2-2k)U_{\max}$:
the coefficient of $U$ is $2 - 2k < 0$ for $k \ge 2$, forcing the adverse
choice $U = +U_{\max}$, while $V$'s coefficient is $+1$, forcing
$V = -V_{\max}$. This is an exact optimisation over the interval box, not a
triangle inequality. The contribution is scaled by $\log 2/4$, as is
$\Delta(k)$ below, both arising through the same
$-\tfrac{\log2}{4}2^{-k}B$ route.

\subsection{The \texorpdfstring{$S_2$}{S2} term: evaluating the Lemma 3.3 error}\label{sec:S2}

Let $\delta_m := A_m - 2^{m+1} + \tfrac{\log2}{4}m^2$ denote the error term of
\cite{GLW}'s Lemma~3.3, and
$\Delta(k) := \sum_i 2^i c_i \delta_{k-i}$ its convolution against the
coefficients --- the error is incurred at every index $k - i$ appearing in
$B$, not once at $k$.

The point of this subsection is that $\delta_m$ can be \emph{evaluated}, not
merely bounded. Setting
$B_j(u) := 2^{-je^{-u}} - 1 + j\log 2\,e^{-u} \ge 0$, one has the exact
decomposition
\begin{equation}\label{eq:delta-decomp}
  \delta_m = -1 - \tfrac{\log 2}{4}m + \sum_{i=0}^{m} 2^i \beta_i(m-i),
  \qquad
  \beta_i(j) := \frac{1}{i!}\int_0^\infty u^i e^{-u} B_j(u)\,du \ge 0 .
\end{equation}
Since every $\beta_i(j) \ge 0$ and $B_j$ is decreasing in $u$, endpoint
bracketing gives certified enclosures of $\delta_m$, with exact
incomplete-gamma panel masses for the $u$-integration.

Bounding $\Delta(k)$ instead by $S^\ast \max_m|\delta_m|$, using \cite{GLW}'s
stated $|\delta_m| \le Cm\log(m+1)$, discards the same sign cancellation
exploited in \S\ref{sec:S3} and costs roughly two orders of magnitude.

In the formalisation the bound actually used is the uniform one,
\begin{equation}\label{eq:Dbound}
  D(m) := (\log 2)\, m \log(m+1) + 1,
  \qquad |\delta_m| \le D(m),
\end{equation}
and \eqref{eq:Dbound} is \emph{proved} rather than assumed: the upper bound
follows because the $\beta_i(j)$ in \eqref{eq:delta-decomp} are non-negative
and $c^\ast(1) < 0$, the lower because
$(\log 2)m\log(m+1) \ge \tfrac{\log2}{4}m$ for $m \ge 1$, and $m = 0$ anchors
at $\delta_0 = -1$, $D(0) = 1$. Thus $S^\ast[2D(k) + D(k+1)]$ in
\eqref{eq:Rbound} is the unsigned route: it discards the cancellation
described above, at the cost of a larger threshold, in exchange for a bound
that requires no $\delta$-table at all.

\begin{remark}
The signed evaluation of \eqref{eq:delta-decomp} is what the certified
$\delta$-table computes, and it is what brings the threshold down; the
formalised route uses \eqref{eq:Dbound} instead. The two are alternatives, not
a chain: \eqref{eq:Dbound} is weaker but self-contained.
\end{remark}

\subsection{The exponentially small term}\label{sec:Eexp}

$E_{\exp}(k)$ collects three contributions: the $s \ge 2$ tail
$2^k \cdot O(1/(k2^k))$, the $2^k \cdot O(3^{-k})$ piece from
$I'_{k,2} = 1/2 + O(3^{-k})$, and the
$\tfrac{\log2}{4}k^2 \cdot O((3/2)^{-k})$ correction from
$B' = d + O((3/2)^{-k})$. Each is bounded by transferring
$\sup|H| \le 101.245$, certified on $|z| = 4.2$, to the relevant inner circles
by the maximum principle. All are below $10^{-180}$ in the range of interest,
so that the constant $3(1.45/k + 10^{-3})$ appearing in \eqref{eq:Rbound} is
enormously conservative.

\subsection{The certified constants}

From \texttt{NumericalFacts}:

\begin{center}
\begin{tabular}{lll}
\toprule
constant & certified bound & r\^ole \\
\midrule
$d$        & $|0.3786950 - d| \le 5.3\times10^{-8}$   & $a_1 = \tfrac14 d\log 2$ \\
$\rho$     & $|-0.1749323 - \rho| \le 6.6\times10^{-8}$ & tail bounds \\
$S^\ast$   & $\le 6.9763$      & \S\ref{sec:S3} \\
$\alpha_+$ & $15.425661$       & $\beta_i$ sweep, \S\ref{sec:S1} \\
$\Gamma$   & $104.5196$        & $\widehat\gamma_i$ sweep, Theorem~6.1 \\
$U_{\max}$ & $0.08436$         & \S\ref{sec:S3}, valid $k \ge 900$ \\
$V_{\max}$ & $1.907462$        & \S\ref{sec:S3}, valid $k \ge 900$ \\
$K_1$      & $2433$            & analytic threshold \\
\bottomrule
\end{tabular}
\end{center}

Here $\alpha_+$ is as in \S\ref{sec:S1}. The constant $\Gamma$ governs the
refined form of the $S_1$ estimate: writing
$\widehat\gamma_i := \max_{s \in [1,2]} |[z^i]G_2^{(2,0)}(s,z)|$ for the
$s$-uniform bound on the \emph{second} $s$-derivatives, set
\[
  \Gamma := \tfrac12 \sum_{i \ge 1} \widehat\gamma_i\, 2^i .
\]
Applying Taylor's theorem with second-order remainder to $g_i$ on $(1,2)$ ---
that is, $g_i(s) = c_i + (s-1)\widehat b_i + R_i(s)$ with
$|R_i(s)| \le (s-1)^2\widehat\gamma_i/2$ --- replaces the first-order estimate
of \S\ref{sec:S1} by one whose remainder is quadratic in $s-1$.

What makes the refinement work is that $\Gamma$ is independent of $k$. This
is the identity $\sum_{l \ge 0} 2^l 3^{-l-1} = 1$: the $l$-th derivative
contributes $2^l$ while the pole at $z = 3$ contributes $3^{-l-1}$, and the
two balance exactly.

The $\widehat\gamma_i$ are evaluated by an interval sweep in Arb using the
closed form $G_2^{(2,0)} = G_2\,(z^2c^2 + zc')$, where $c = c(s,z)$ is as in
\S\ref{sec:S1}; that closed form is proved rather than assumed. At the pinned
settings --- primes to $P = 10^5$, $800$ boxes, $300$ bits --- the sweep gives
$\Gamma = 104.5196$, as a head contribution $203.118$ and a tail $5.921$
against $\sup|G_2^{(2,0)}|\le 19881.2$ on $|z| = 2.9$.

The refined estimate is what lowers the sweep threshold to $K_0 = 78$, and so
what lets \texttt{sweep\_mid} begin at $k = 79$ rather than higher; the
unrefined first-order route of \S\ref{sec:S1} would leave a gap between the
sweeps.

\begin{remark}\label{rem:two-routes}
Two routes through \S\ref{sec:S2} are available, and we have taken the weaker
one. The uniform bound \eqref{eq:Dbound} is self-contained and proved; it
costs a larger analytic threshold, which the certified sweeps then absorb. The
$\delta$-table route evaluates \eqref{eq:delta-decomp} and exploits the sign
cancellation, giving a much smaller threshold at the price of a substantial
certified computation. Since the sweeps cover the intervening range either
way, the theorem does not depend on the choice, and we prefer the route that
rests on less.
\end{remark}

From \texttt{NumericalFacts}, with the values from the current certified run:

%=======================================================================
\section{Assembly and the certified ranges}\label{sec:full-assembly}
%=======================================================================

The proof of Theorem~\ref{thm:main-full} splits $[6,\infty)$ into four ranges.
The first three are certified computations; the fourth is the analytic
argument of Proposition~\ref{prop:crit-ge-phi}.

\begin{center}
\begin{tabular}{lll}
\toprule
range & source & field \\
\midrule
$6 \le k \le 199$    & the prior certified $f_k$ computation      & \texttt{sweep\_low} \\
$200 \le k \le 899$  & refined criterion sweep ($K_0 = 78$)      & \texttt{sweep\_mid} \\
$900 \le k \le 2439$ & extended $\delta$-table sweep             & \texttt{sweep\_high} \\
$k \ge 2440$         & derivative domination, proved here        & \texttt{phi\_at\_k1}, \texttt{phi\_mono} \\
\bottomrule
\end{tabular}
\end{center}

\noindent
The ranges overlap pairwise, so their union is $[6,\infty)$ and
Theorem~\ref{thm:main-full} follows.

For the analytic range, Proposition~\ref{prop:crit-ge-phi} reduces matters to
positivity of $\varphi$. The certified scan gives
$\varphi(2440) = 89 \pm 0.35 > 0$, and $\varphi(m) \le \varphi(m+1)$ for every
$m \ge 2440$; induction gives $\varphi(k) > 0$, hence $f_k < f_{k+1}$, for all
$k \ge 2440$.

\begin{remark}
The monotonicity of $\varphi$ above $2440$ is a statement about
\eqref{eq:Rbound}: the leading term $a_1 k^2$ grows quadratically while $R(k)$
grows linearly in $k$ apart from the $\delta$-table contribution, which is
controlled by the same certified table.
\end{remark}

%=======================================================================
\section{A general principle}\label{sec:general}
%=======================================================================

\begin{conjecture}\label{conj:general}
Let $\mathcal{P}$ be a set of primes with least element $q$, and let
$f_k^{\mathcal{P}} = \sum 1/(n \log n)$ over those $n$ with $\Om(n) = k$ and
all prime factors in $\mathcal{P}$. Then the $k^2$ term is present --- and
hence the sequence is eventually increasing --- if and only if $q = 2$.
\end{conjecture}

The heuristic is Proposition~\ref{prop:analyticity}. The regularised product
$H^{\mathcal{P}}(s,\cdot)$ is analytic in $|z| < q^{s}$, and the evaluation
point relevant to the $2^{-k}$ scale is $z = 2$; so $z = 2$ lies on the
boundary at $s = 1$ precisely when $q = 2$, and strictly inside otherwise. On
the boundary the $k^2$ factor appears and competes with $2^{-k}$; strictly
inside it does not.

Theorem~\ref{thm:main-full} is the case $\mathcal{P} = \{\text{all primes}\}$,
$q = 2$, and Theorem~\ref{thm:main-odd} is the case
$\mathcal{P} = \{p \ge 3\}$, $q = 3$. For $q \ge 3$ the argument of
\S\ref{sec:odd} applies with only the constants changing: the limit is the
omitted Euler factors evaluated at $z = s = 1$, the coefficient of $2^{-k}$ is
$H^{\mathcal{P}}(1,2)\,(\partial_s \log H^{\mathcal{P}}(1,2) + 2\gamma)$, and
the remainder is controlled by Cauchy's estimate after factoring out the pole
at $z = q^{s}$.

Two caveats keep this a conjecture rather than a theorem. The first concerns
the ``if'' direction for general $\mathcal{P}$: the argument of
\S\ref{sec:odd} needs the regularised product to converge and its coefficients
to be extractable, and while $\sum_{p \in \mathcal{P}} p^{-2s}$ converges for
any $\mathcal{P}$, the identification of the limit and of the coefficient of
$2^{-k}$ uses the full Euler product for $\zeta$. For $\mathcal{P}$ the set of
all primes above a bound, or a union of arithmetic progressions of primes, the
missing factors are controlled and the argument should go through unchanged;
for arbitrary $\mathcal{P}$ we do not claim it.

The second concerns the ``only if'' direction. That $2 \in \mathcal{P}$ places
$z = 2$ on the boundary is immediate, but that the $k^2$ term is then
genuinely present --- rather than vanishing by some cancellation --- is a
statement we have verified only for $\mathcal{P}$ the set of all primes, where
it is \cite{GLW}.

\begin{remark}
The two theorems of this paper are the two extreme cases: removing the single
prime $2$ already moves the radius past the evaluation point, and no smaller
change does. In that sense the conjecture asserts that nothing between them
occurs.
\end{remark}

%=======================================================================
\section{What is proved and what is assumed}\label{sec:assumed}
%=======================================================================

We separate three things: what is proved in Lean, what is computed in ball
arithmetic and taken as a hypothesis, and what is inherited from
\cite{GLW}.

\begin{itemize}
  \item The analytic reductions of \S\S\ref{sec:mechanism}--\ref{sec:odd} and
        \S\S\ref{sec:full-criterion}--\ref{sec:full-assembly} are formalised in
        Lean~4 with no \texttt{sorry}, resting only on \texttt{propext},
        \texttt{Classical.choice} and \texttt{Quot.sound}.
  \item The numerical inputs are ball-arithmetic computations, taken as
        hypotheses rather than proved within the formalisation. For
        Theorem~\ref{thm:main-full} these are collected in a structure
        \texttt{NumericalFacts}: the enclosures of $d$ and $\rho$, the bounds
        $S^\ast \le 6.9763$, $\alpha_+ \le 15.4257$ and $\Gamma \le 104.52$,
        the aggregate source bound and the two $\varphi$ facts at the
        analytic threshold, and the three sweeps
        $[6,199]$, $[79,899]$, $[900,2439]$. For
        Theorem~\ref{thm:main-odd} the corresponding structure carries the
        enclosure of $\Sodd$, the bound on $M_R$, the expansion
        \eqref{eq:odd-expansion} with its remainder bound, and the sweep at
        small $k$.
  \item Consequently both theorems are proved \emph{conditionally on} those
        computations, in the same sense and to the same extent. What the
        formalisation establishes is that no gap remains in the reduction:
        given the certified inputs, the analytic argument is complete and
        machine-checked.
  \item A second form of the main theorem replaces the aggregate source
        bound by the quantities the certified package actually computes ---
        the $\beta_i$ sweep, $M_0$, the $S^\ast$ and $\alpha_+$ partial-sum
        bounds, the $U$/$V$ and $E_{\exp}$ sources --- and derives the source
        bound from them, so that \S\S\ref{sec:S1}--\ref{sec:Eexp} are
        load-bearing rather than merely stated.
  \item Two facts absent from Mathlib were proved in the course of the odd
        formalisation and may be of independent use:
        $\lim_{s \to 1^+}(\zeta(s) - 1/(s-1)) = \gamma$, and that
        $\zeta(s) - 1/(s-1)$ extends analytically across $s = 1$.
\end{itemize}

%=======================================================================
\appendix
\section{Certified computations}\label{app:numerics}
%=======================================================================

All interval computations use Arb/FLINT through python-flint. The Stieltjes
constants entering the representation of $\zeta(s)(s-1)$ in \S\ref{sec:S1},
and the $\zeta$ evaluations used in boundary checks, are FLINT's own certified
implementations. The $\delta$-table of \S\ref{sec:S2} was cross-validated
against an independently written high-precision implementation of the original
integral definition of $A_m$, at $40$--$80$ decimal digits, agreeing to six
or more significant figures at each test point.

The scripts are \texttt{numerics/} for Theorem~\ref{thm:main-full} and
\texttt{odd-proof/numerics/} for Theorem~\ref{thm:main-odd}. The two sweeps behind
$\alpha_+$ and $\Gamma$ use $1000$ and $600$ boxes respectively in
$s \in [1,2]$.

For the odd case, $M_R$ is computed over $40$ boxes in $s \in [1,2]$ and $100$
boxes in $\theta$ on the circle, with primes to $5 \times 10^4$ and an
explicit tail bound. Subdividing in $s$ is essential rather than cosmetic: a
single box over $[1,2]$ gives only $M_R \le 739605$.

The five constants of \eqref{eq:odd-calibrated} beyond $S$ are certified as
follows.
\begin{center}
\begin{tabular}{lll}
\toprule
constant & bound & method \\
\midrule
$S_0$   & $\le 44.6$   & first-order pole; recursion (head), Cauchy (tail) \\
$S_1$   & $\le 453.1$  & second-order pole, recursion \\
$C_N$   & $\le 364$    & odd Euler-product recursion $+$ Cauchy tail \\
$N_U$   & $\le 24.68$  & exact head to $k = 400$, explicit envelope beyond \\
$C_E$   & $\le 5.48910$ & Taylor model in $\varepsilon = s-1$ \\
\bottomrule
\end{tabular}
\end{center}
\noindent
$C_E$ is certified over the whole range $s \in (1,\,s(3)]$ rather than at a
sample of points. The four constants making it up are ratios in which the pole
of $\zeta$ at $s = 1$ cancels, so interval evaluation on boxes of $s$ loses the
cancellation and diverges however narrow the boxes are taken; expanding $\zeta$
as a power series about $1 + E/2$ makes the cancellation symbolic, and the
resulting series --- whose coefficients decay --- evaluates tightly on the
whole disc in one step. The scripts are \texttt{s0\_recursion.py},
\texttt{s1\_recursion.py}, \texttt{cn\_recursion.py}, \texttt{nearu\_bound.py}
and \texttt{ce\_taylor.py}.

The two sweeps behind $\alpha_+$ and $\Gamma$ are run at pinned settings:
primes to $P = 10^5$ and $300$ bits of precision in both, with $1600$ boxes in
$s$ for $\alpha_+$ and $800$ for $\Gamma$. The $f_k$ certificate records its
own parameters --- $2194$ panels of width $0.8$, $52$ Gauss--Legendre points
per panel, $1400$ bits, Bernstein-ellipse parameter $1.45$ --- together with a
SHA-256 hash of the payload excluding the hash field itself.

The quadrature error is controlled by \cite[Thm.~19.3]{ATAP}: if $f$ is
analytic on $[-1,1]$ and continues analytically to the open Bernstein ellipse
$E_\rho$ --- the ellipse with foci $\pm 1$ whose semimajor and semiminor axes
sum to $\rho$ --- with $|f| \le M$ there, then $(n+1)$-point Gauss quadrature
satisfies, for $n \ge 2$,
\begin{equation}\label{eq:trefethen}
  \Bigl|\int_{-1}^{1} f - \widehat{I}\,\Bigr|
    \;\le\; \frac{144}{35}\,\frac{M\rho^{-2n}}{\rho^2 - 1}.
\end{equation}
The panel decomposition applies \eqref{eq:trefethen} on each of the $2194$
panels, with $\rho$ determined by the distance from the panel to the nearest
singularity of the integrand.

The certified run records python-flint $0.9.0$ under Python $3.12.3$.

In the implementation the per-panel bound is applied in the equivalent form
\[
  10 \cdot \frac{64}{15}\,
    \frac{\rho^{-2N}}{1 - \rho^{-2}},
  \qquad N = 52,
\]
which coincides identically with \eqref{eq:trefethen} at $n = N - 1 = 51$,
since $\rho^{-2N}/(1-\rho^{-2}) = \rho^{-2(N-1)}/(\rho^2-1)$; the constant
$64/15$ is used in place of $144/35$, which is larger by $3.7\%$ and hence
conservative. With panel width $H = 0.8$ and ellipse parameter
$\delta = 1.45$ one has $\rho = \delta' + \sqrt{1 + \delta'^2}$ for
$\delta' = 2\delta/H$, giving $\rho \approx 7.385$ and a per-panel
discretisation error below $10^{-88}$ --- far smaller than the arithmetic and
truncation errors it accompanies.

%=======================================================================
\begin{thebibliography}{9}

\bibitem{BanksMartin}
W.~D. Banks and G.~Martin,
\emph{Optimal primitive sets with restricted primes},
Integers \textbf{13} (2013), \#A69, 10 pp.; arXiv:1301.0948.

\bibitem{GLW}
O.~Gorodetsky, J.~D. Lichtman and M.~W. Wong,
\emph{On Erd\H{o}s sums of almost primes},
C.~R. Math. Acad. Sci. Paris \textbf{362} (2024), 1571--1580;
arXiv:2303.08277.

\bibitem{Erdos}
P.~Erd\H{o}s,
\emph{Note on sequences of integers no one of which is divisible by any
other},
J.~London Math. Soc. \textbf{10} (1935), 126--128.

\bibitem{ATAP}
L.~N. Trefethen,
\emph{Approximation Theory and Approximation Practice},
SIAM, Philadelphia, 2013.

\bibitem{Zhang}
Z.~Zhang,
\emph{On a problem of Erd\H{o}s concerning primitive sequences},
Math. Comp. \textbf{60} (1993), 827--834.

\bibitem{BKK}
J.~Bayless, P.~Kinlaw and D.~Klyve,
\emph{Sums over primitive sets with a fixed number of prime factors},
Math. Comp. \textbf{88} (2019), 3063--3077.

\end{thebibliography}

\end{document}
